\documentclass[10pt,a4paper]{amsart}
\usepackage[margin=19mm]{geometry}
\usepackage{amsmath,amssymb,mathtools}
\usepackage[scr=rsfs,cal=euler]{mathalfa}
\usepackage{xfrac}
\usepackage[unicode,hidelinks,bookmarks=false]{hyperref}
\newcommand{\ie}{i.e.\ }
\newcommand{\cf}{cf.\ }
\newcommand{\resp}{respectively\ }
\newcommand{\Subsection}{\subsection}
\providecommand{\nobreakdash}{\nobreak\hbox{-}}
\setlength{\emergencystretch}{2em}
\multlinegap0pt
\usepackage{bm}
\usepackage{bbm}
\usepackage{dsfont}
\usepackage{xspace}
\usepackage{cancel}
\usepackage{mathtools}
\usepackage{amsthm}
\makeatletter
\@ifundefined{thm}{\newtheorem{thm}{Theorem}}{}
\makeatother
\DeclareMathOperator{\Spec}{Spec}
\DeclareMathOperator{\Sym}{Sym}
\newcommand{\Pic}{\operatorname{Pic}}
\newcommand{\dendegs}{\delta}
\newcommand{\nn}{\mathbb{N}}
\title[Density of algebraic points on products of curves]{Density of algebraic points on products of curves}
\author{Jennifer Berg, Yu Fu, Evangelia Gazaki, Morena Porzio, James Rawson, Isabel Vogt}
\date{Source: arXiv:2507.00860v1 (CC BY 4.0)}
\begin{document}
\maketitle

\noindent\emph{Source and licence.} Density of algebraic points on products of curves. Version arXiv:2507.00860v1 (CC BY 4.0), distributed under the Creative Commons Attribution 4.0 International licence, as recorded in the arXiv licence metadata for this exact version and shown on the abstract page https://arxiv.org/abs/2507.00860v1. Full rights notice: licenses/arxiv-2507.00860-CC-BY-4.0.txt.\par\smallskip

\noindent\textit{Editorial scope.} Counted unit: the Thm whose source label is \texttt{thm:Jac of genus 2} in the article (source0.tex lines 1349--1352, source label \texttt{thm:Jac of genus 2}), delivered together with the article's own complete original proof of it (source0.tex lines 1353--1411, one \texttt{proof} environment of 59 lines). No mathematics is written, repaired or re-derived: statement and proof are the article's own text line for line. Required context retained uncounted: none. Every definition, display and label of those lines is retained unchanged. Omitted from this package and not redistributed: 1--1348, 1412--1734. Nothing inside a retained excerpt is omitted or altered: every retained body is delivered exactly as the article's lines are written, so no omission receipt of any kind accompanies this package. Citations: the retained excerpts carry no citation command at all, so no bibliography entry of the article is transcribed and no reference is redistributed. Reference adaptation: none was needed and none was made; every reference inside the delivered unit resolves inside it, so no unresolved reference or citation is produced. Printed slips kept literal: no printed word, symbol, hypothesis or number of the counted statement or of its proof was altered. Layout and class: the article's own class is \texttt{amsart} with packages including babel, amsfonts, amsthm, amsmath, amssymb, enumerate, mathtools, tikz, tikz-cd, xy, graphicx, adjustbox, extpfeil, comment; the wrapper is the lane's pinned \texttt{amsart} editorial wrapper at 10pt on A4, which restates the theorem-like environments the retained text needs and adds the article's own macro definitions that the retained text needs (\texttt{\textbackslash Pic}, \texttt{\textbackslash Spec}, \texttt{\textbackslash Sym}, \texttt{\textbackslash dendegs}, \texttt{\textbackslash nn}), with extra wrapper packages (bm, bbm, dsfont, xspace, cancel, mathtools). The delivered numbering is the unit's own. Assets: no retained span contains a figure environment, a graphics-inclusion command, a PDF-page inclusion, a table or a TikZ picture; no image file is copied into this package, the package declares no asset dependency and no image file is redistributed with it. Licence: version arXiv:2507.00860v1 of this article is distributed under the Creative Commons Attribution 4.0 International licence, as recorded in the arXiv licence metadata for this exact version and shown on the abstract page https://arxiv.org/abs/2507.00860v1. A full rights notice is delivered at licenses/arxiv-2507.00860-CC-BY-4.0.txt. Characters: every retained excerpt is pure ASCII, so the delivered engine, XeLaTeX with its default Latin Modern text font, reports no missing character.
\begin{thm}\label{thm:Jac of genus 2} 
Let $C$ be a nice genus $2$ curve over $k$. 
Then $\nn_{\geq 2}\subset\dendegs(\Pic^0_C/k)$. 
\end{thm}

\begin{proof}
    If $\Pic^0_C$ contains a curve $D$ of genus at least 2, then the images of the curve under $[n]$ are dense. 
    In particular, the density degree set of $\Pic^0_C$ contains the density degree set of $D$. 
    First, consider the map 
    \[
        C\xhookrightarrow{\Delta_C} C\times C \to \Sym^2_C \overset{AJ_C}{\to} \Pic^2_C \overset{\otimes \omega_C^{-1}}{\to} \Pic^0_C.
    \]
    The image of the composite is birational to $C$ (only the Weierstrass points get identified), and therefore $\dendegs(C/k)\subseteq \dendegs(\Pic^0_C/k)$. Since $C$ is hyperelliptic, this gives $2\nn\subset \dendegs(\Pic^0_C/k)$.
    
    We now show that there exists a curve $D$ in $\Pic^0_C$ of arithmetic genus 4 with dense degree 3 points, such that its normalization $D^\nu$ has a rational point (giving index 1).
    We may write $C : y^2 = g(x) = b_6 x^6 + b_4 x^4 + b_3 x^3 + b_2 x^2 + b_1 x + b_0$ with at least one of $b_3$ and $b_1$ non-zero. 
    If not, applying the M\"obius transformation $\frac{1}{tx + 1} - s$ gives another model, $C'$, for $C$, where $t \in k^\times$, and $s$ is chosen such that the new model has $b_5 = 0$. If for the new model $b_3 = b_1 = 0$ as well, there is an automorphism of $C'$ given by $(x, y) \mapsto (-x, y)$. Composing this with the change of model gives an automorphism of $C$, mapping the $x$-coordinate to $-\frac{(2st + 1)x + 2s}{(2st^2 + 2t)x + (2st + 1)}$. These automorphisms will be distinct for different $t$, which implies that for all but finitely many $t$, the new model will satisfy this assumption.
    
    Define $X_\sigma = C \times_{\sigma} C$, where $\mathbb{P}^1_k \xrightarrow{\sigma} \mathbb{P}^1_k $ is the automorphism sending $x\mapsto -x$ and $C\to \mathbb{P}^1_k$ is the usual $2:1$ cover. 
    Notice that the $S_2$-action on $C\times C$ (swapping the coordinates) restricts to the closed subscheme $X_\sigma$, which means that $X_\sigma$ descends to a closed subscheme $Z_\sigma$ in $\Sym^2_C$. 
    If $\{\infty_1,\infty_2\}\subseteq C$ are the two $\overline{k}$-points above $\infty\in \mathbb{P}^1_k$, then the orbit $\{(\infty_1),(\infty_2)\}$ descends to a unique $k$-point on $Z_\sigma$, and hence index$(Z_\sigma)=1$.
    Over the affine chart
    \[
    \Spec R = \Spec \left( k[x ,y_1,y_2]/ (y_1^2- g( x ), y_2^2- g(\sigma x )) \right)
    \]
    of $X_\sigma$, the $S_2$-action on $X_\sigma$ can be written as $(x, y_1, y_2) \mapsto ( \sigma x, y_2, y_1)$. 
    Consider the ring of invariants
    \[
    R^{S_2} = \left( k[x ,y_1,y_2]/ (y_1^2- g( x ), y_2^2- g(\sigma x )) \right)^{S_2},
    \]
    so that $\Spec(R^{S_2})$ is an affine chart of $Z_\sigma$.
    Notice that $u=x^2, w_1= y_1+y_2,  w_2 = y_1y_2$ are all invariants.
    Moreover we have the element $v = \frac{y_1-y_2}{x}\in R_x^{S_2}$ as well.
   In $R_x^{S_2}$ we can also write $w_2$ as $\frac{1}{4}\left(w_1^2 - uv^2\right)$.
    Let us check that $u,v,w_1$ generate the ring of invariants $R_x^{S_2}$. 
    Pick 
    \[
    p(x,y_1,y_2) = \sum_{i\in \mathbb{Z}} \alpha_i x^i + \sum_{i\in \mathbb{Z}} \beta_i x^iy_1 +\sum_{i\in \mathbb{Z}} \gamma_i x^i y_2  + \sum_{i\in \mathbb{Z}} \delta_i x^i y_1y_2 \in R_x^{S_2}
    \]
    and apply the $S_2$ action.
    We then get
    \[
    p(x,y_1,y_2) = \sum_{i\in \mathbb{Z}} (-1)^i \alpha_i x^i + \sum_{i\in \mathbb{Z}} (-1)^i \beta_i x^iy_2 +\sum_{i\in \mathbb{Z}} (-1)^i \gamma_i x^i y_1  + \sum_{i\in \mathbb{Z}} (-1)^i \delta_i x^i y_1y_2,
    \]
    which implies that $\delta_{odd}=\alpha_{odd}=0$ and that $(-1)^i\beta_i=\gamma_i$.
    This means that $p(x,y_1,y_2)$ can be written as a polynomial in $u,v,w_1$.
    Notice that these elements satisfy the following relations:
    \[
    vw_1 = \frac{(y_1-y_2)(y_1+y_2)}{x} = \frac{y_1^2 - y_2^2}{x}  = 2(b_3x^2+b_1) = 2(b_3u+b_1),
    \]
    \[
    \frac{1}{2}(w_1^2+ uv^2) = w_1^2 - 2w_2 = 2b_6 u^3 + 2b_4 u^2 + 2b_2 u + 2b_0.
    \]
    
    This realizes the open $\Spec( R_x^{S_2} )$ of $Z_\sigma$ as a closed of 
    \[
        U = \Spec \left( k[u,v,w_1]/ (vw_1-(2b_3u+2b_1), w_1^2+uv^2 - (4b_6u^3 + 4b_4u^2 + 4b_2u + 4b_0)) \right).
    \]
    Since $U$ is integral and $Z_\sigma$ has the same dimension as $U$, then $U=\Spec( R_x^{S_2})$.
    In particular $Z_\sigma$ is birational to the intersection of a quadric and a cubic and  $D=\overline{U}\to Z_\sigma$ is a birational morphism with $D$ a curve with arithmetic genus 4.
    If $b_3 \neq 0$, then $u = \frac{vw_1 - 2b_1}{2b_3}$ equips the quadric with a ruling over the base field, and thus $X_\sigma$ with degree 3 morphisms to $\mathbb{P}^1$ (namely $v$ and $w_1$). Otherwise, $vw_1 = b_1$, and eliminating $v$ from the second equation shows $b_1^2 u + w_1^4 = 2(b_6 u^3 + b_4 u^2 + b_2 u + b_0)w_1^2$. This also has a degree 3 morphism given by $w_1$. 

    Pushing forward $Z$ under the natural maps $\Sym^2_C \to \Pic^2_C \to \Pic^0_C$ gives a curve in $\Pic^0_C$. 
\end{proof}

\end{document}
